\SourceLeaf{233}%
dieses Gegenstands für Figuren in größtmöglicher Allgemeinheit müssen wir uns jedoch für eine andere Gelegenheit vorbehalten.\par\endgroup

\Article{7}
Suchen wir nun eine Formel für das Krümmungsmaß in einem beliebigen Punkt einer gekrümmten Fläche. Bezeichnet \(\mathrm{d}\sigma\) den Flächeninhalt eines Elements dieser Fläche, so ist \(Z\mathrm{d}\sigma\) der Flächeninhalt seiner Projektion auf die Ebene der Koordinaten \(x,y\). Ist entsprechend \(\mathrm{d}\Sigma\) der Flächeninhalt des zugehörigen Elements auf der Kugeloberfläche, so ist \(Z\mathrm{d}\Sigma\) der Flächeninhalt seiner Projektion auf dieselbe Ebene. Das positive oder negative Vorzeichen von \(Z\) zeigt an, ob die Projektion ähnlich oder entgegengesetzt zum projizierten Element liegt. Offenbar stehen die Projektionen also hinsichtlich der Größe im selben Verhältnis und hinsichtlich der Lage in derselben Beziehung zueinander wie die Elemente selbst. Betrachten wir nun ein Dreieckselement auf der gekrümmten Fläche und nehmen an, die Koordinaten der drei Punkte, die seine Projektion bilden, seien
\[
\begin{array}{ll}
x,&y\\
x+\mathrm{d}x,&y+\mathrm{d}y\\
x+\delta x,&y+\delta y
\end{array}
\]
Der doppelte Flächeninhalt dieses Dreiecks wird durch die Formel
\[
\mathrm{d}x\mathord{.}\delta y-\mathrm{d}y\mathord{.}\delta x
\]
ausgedrückt, und zwar positiv oder negativ, je nachdem, ob die Seite vom ersten zum dritten Punkt in Bezug auf die Seite vom ersten zum zweiten Punkt ebenso liegt wie die Koordinatenachse \(y\) in Bezug auf die Koordinatenachse \(x\) oder ob sie entgegengesetzt liegt.

Entsprechend gilt für das zugehörige Element auf der Kugeloberfläche Folgendes. Haben die drei Punkte, die seine Projektion bilden, mit dem Kugelmittelpunkt als Ursprung die Koordinaten
\[
\begin{array}{ll}
X,&Y\\
X+\mathrm{d}X,&Y+\mathrm{d}Y\\
X+\delta X,&Y+\delta Y
\end{array}
\]
so ergibt sich für den doppelten Flächeninhalt dieser Projektion
\[
\mathrm{d}X\mathord{.}\delta Y-\mathrm{d}Y\mathord{.}\delta X
\]
Für das Vorzeichen dieses Ausdrucks gilt dasselbe wie oben. Somit ist an dieser Stelle der gekrümmten Fläche %
